%% kapa: exam=1 round=30 year=2019 period=2 form=A subject=accounting subject_ko=회계학 seq=20 number=60
%% source: Q-Net 감정평가사 기출문제 (cst00309) artlSeq=5205205 file=extracted/1008/5205205/2019년 제30회 감정평가사 1차시험 문제지/2019년 제30회 감정평가사 1차시험 문제지/2019년도 제30회 감정평가사 1차 2교시 A형.hwp
%% conversion: kordoc 4.24.51 -> tools/gen1.py (auto); review=proofread-sonnet; flags=none
\begin{question}{60}
(주)감평은 20x1년 1월 1일 액면금액 ￦1,000,000(만기 3년, 표시이자율 연 6 \%, 매년 말 이자지급)의 사채를 발행하였으며, 사채의 발행 당시 유효이자율은 연 8 \%이었다. (주)감평은 20x2년 6월 30일 사채를 조기상환하였다. 조기상환시 발생한 사채상환손실은 ￦32,000이다. (주)감평이 유효이자율법을 적용할 때, 상환일까지의 경과이자를 포함한 사채조기상환금액은? (단, 이자비용은 월할계산하고, 계산금액은 소수점 첫째자리에서 반올림하며, 단수차이로 인한 오차가 있으면 가장 근사치를 선택한다.) \par
{\small\centering
\begin{tabular}{|c|c|c|c|c|}
\hline
\multirow{2}{*}{\makecell{기간}} & \multicolumn{2}{|c|}{\makecell{단일금액 ￦1의 현재가치}} & \multicolumn{2}{|c|}{\makecell{정상연금 ￦1의 현재가치}} \\
\cline{2-5}
 & \makecell{6 \%} & \makecell{8 \%} & \makecell{6 \%} & \makecell{8 \%} \\
\hline
\makecell{1} & \makecell{0.9434} & \makecell{0.9259} & \makecell{0.9434} & \makecell{0.9259} \\
\hline
\makecell{2} & \makecell{0.8900} & \makecell{0.8574} & \makecell{1.8334} & \makecell{1.7833} \\
\hline
\makecell{3} & \makecell{0.8396} & \makecell{0.7938} & \makecell{2.6730} & \makecell{2.5771} \\
\hline
\end{tabular}\par}

\begin{choices}
  \item ￦970,872
  \item ￦996,300
  \item ￦1,004,872
  \item ￦1,034,872
  \item ￦1,073,444
\end{choices}
\end{question}
\answer{4}
